\chapter{Degree of Maps}
\label{ch:viii03-degree}

Homotopy theory becomes effective when a deformation-invariant quantity can be
computed.  Degree is the first such integer-valued invariant in this volume.
For a map between oriented spaces of the same dimension, degree records the
signed number of times the domain covers the target.  It detects maps that
cannot be deformed into one another and provides immediate applications to
spheres and antipodal maps.

\section{Orientation and the signed-counting idea}

For a smooth map
\[
f:M^n\to N^n
\]
between oriented \(n\)-manifolds, a point \(x\in M\) is regular if the derivative
\[
df_x:T_xM\to T_{f(x)}N
\]
is an isomorphism.  At such a point, \(df_x\) either preserves or reverses
orientation.

\begin{definition}[Local sign]
\label{def:viii03-local-sign}
At a regular point \(x\), define
\[
\operatorname{sgn}_x(f)=
\begin{cases}
+1,&df_x\text{ preserves orientation},\\
-1,&df_x\text{ reverses orientation}.
\end{cases}
\]
\end{definition}

\section{Degree by a regular value}

\begin{definition}[Degree for a smooth map]
\label{def:viii03-degree}
Let \(M,N\) be connected compact oriented \(n\)-manifolds without boundary,
and let \(f:M\to N\) be smooth.  If \(y\in N\) is a regular value, define
\[
\deg(f)
=
\sum_{x\in f^{-1}(y)}
\operatorname{sgn}_x(f).
\]
\end{definition}

Compactness makes the regular fiber finite.

\begin{theorem}[Independence of regular value]
\label{thm:viii03-regular-independence}
The signed sum above is independent of the chosen regular value \(y\).
\end{theorem}

This theorem is the key reason degree is global rather than merely local.
A full proof can be given by oriented cobordism of inverse images of a path
between regular values.  Later homology theory will provide a shorter algebraic
formulation.

\section{Continuous maps}

\begin{theorem}[Continuous degree]
\label{thm:viii03-continuous}
Every continuous map \(f:M\to N\) between connected closed oriented
\(n\)-manifolds has an integer degree.  It agrees with the signed regular-value
count for smooth maps and is invariant under homotopy.
\end{theorem}

One construction approximates a continuous map by a smooth map and proves that
different sufficiently close approximations are homotopic.  Later, after
homology is available, degree will be characterized by the induced map on the
top-dimensional fundamental class.

\section{Normalization}

\begin{proposition}[Identity and constant map]
\label{prop:viii03-normalization}
For a connected oriented closed \(n\)-manifold,
\[
\deg(\operatorname{id})=1.
\]
If \(n>0\), every constant map has degree \(0\).
\end{proposition}

\begin{proof}
The identity has one preimage of a regular value and preserves orientation.
A constant map misses every point other than its image, so choosing a regular
value outside that image gives an empty signed sum.
\end{proof}

\section{Homotopy invariance}

\begin{theorem}[Homotopy invariance of degree]
\label{thm:viii03-homotopy}
If
\[
f\simeq g:M\to N,
\]
then
\[
\deg(f)=\deg(g).
\]
\end{theorem}

\begin{corollary}[Nonzero degree obstructs null-homotopy]
\label{cor:viii03-nonzero}
If \(n>0\) and \(\deg(f)\ne0\), then \(f\) is not null-homotopic.
\end{corollary}

\begin{proof}
A null-homotopic map is homotopic to a constant map, whose degree is \(0\).
\end{proof}

\section{Multiplicativity}

\begin{theorem}[Degree of a composite]
\label{thm:viii03-multiplicative}
For maps of connected closed oriented \(n\)-manifolds
\[
M\xrightarrow{f}N\xrightarrow{g}P,
\]
one has
\[
\deg(g\circ f)=\deg(g)\deg(f).
\]
\end{theorem}

\begin{corollary}[Degree of a homotopy equivalence]
\label{cor:viii03-equivalence}
If \(f:M\to N\) is a homotopy equivalence between connected closed oriented
\(n\)-manifolds, then
\[
\deg(f)=\pm1.
\]
\end{corollary}

\begin{proof}
For a homotopy inverse \(g\),
\[
1=\deg(\operatorname{id}_M)
=\deg(g\circ f)
=\deg(g)\deg(f).
\]
The only units in \(\mathbb Z\) are \(\pm1\).
\end{proof}

\section{Degree on the circle}

Identify
\[
S^1=\{z\in\mathbb C:|z|=1\}.
\]

\begin{theorem}[Power map]
\label{thm:viii03-power}
For
\[
p_k:S^1\to S^1,
\qquad
p_k(z)=z^k,
\]
with \(k\in\mathbb Z\),
\[
\deg(p_k)=k.
\]
\end{theorem}

For \(k>0\), a regular value has \(k\) preimages, each with positive local sign.
For \(k<0\), complex conjugation reverses orientation and contributes the
negative sign.

\begin{example}
\[
\deg(z\mapsto z^5)=5,
\qquad
\deg(z\mapsto z^{-3})=-3.
\]
\end{example}

\section{Degree classifies circle maps up to homotopy}

\begin{theorem}[Classification of maps \(S^1\to S^1\)]
\label{thm:viii03-circle-classification}
Two continuous maps
\[
f,g:S^1\to S^1
\]
are homotopic if and only if
\[
\deg(f)=\deg(g).
\]
\end{theorem}

This is an early instance of a complete homotopy classification by an algebraic
invariant.  Later the fundamental group identifies the same integer as a
winding number.

\section{Orientation-reversing maps}

\begin{proposition}[Reflection]
\label{prop:viii03-reflection}
A reflection of \(S^n\) across a hyperplane has degree \(-1\).
\end{proposition}

\begin{proposition}[Orientation-preserving homeomorphism]
\label{prop:viii03-homeomorphism}
An orientation-preserving homeomorphism of an oriented closed manifold has
degree \(+1\); an orientation-reversing one has degree \(-1\).
\end{proposition}

\section{The antipodal map}

Let
\[
a:S^n\to S^n,
\qquad
a(x)=-x.
\]

\begin{theorem}[Degree of the antipodal map]
\label{thm:viii03-antipodal}
The antipodal map on \(S^n\) has degree
\[
\deg(a)=(-1)^{n+1}.
\]
\end{theorem}

\begin{proof}
The antipodal map is the restriction of \(-I:\mathbb R^{n+1}\to\mathbb R^{n+1}\).
Its determinant is
\[
\det(-I)=(-1)^{n+1}.
\]
The induced orientation sign on the sphere agrees with this value.
\end{proof}

Thus the antipodal map has degree \(+1\) for odd-dimensional spheres and
degree \(-1\) for even-dimensional spheres.

\section{No-retraction preview}

Degree can obstruct retractions.

\begin{theorem}[No retraction of a ball onto its boundary]
\label{thm:viii03-no-retraction}
For \(n\ge1\), there is no retraction
\[
r:D^{n+1}\to S^n.
\]
\end{theorem}

\begin{proof}
Suppose \(r\) existed and let \(i:S^n\hookrightarrow D^{n+1}\) be inclusion.
Then
\[
r\circ i=\operatorname{id}_{S^n}.
\]
Because \(D^{n+1}\) is contractible, \(i\) is null-homotopic as a map into the
disk, so \(r\circ i\) is null-homotopic by postcomposition.  But the identity
of \(S^n\) has degree \(1\), while a null-homotopic map has degree \(0\), a
contradiction.
\end{proof}

\section{Surjectivity consequence}

\begin{proposition}[Nonzero degree implies surjective]
\label{prop:viii03-surjective}
If \(f:M\to N\) is a continuous map between connected closed oriented
\(n\)-manifolds and
\[
\deg(f)\ne0,
\]
then \(f\) is surjective.
\end{proposition}

\begin{proof}
If \(f\) missed a point \(y\in N\), then in the smooth regular-value picture a
nearby regular value outside the image would have empty preimage and hence
degree \(0\).  Equivalently, a map missing a point factors through a punctured
target, which has no top-dimensional degree contribution.
\end{proof}

\section{Degree and homotopy equivalence of spheres}

\begin{corollary}
A self-map \(f:S^n\to S^n\) of degree other than \(\pm1\) cannot be a homotopy
equivalence.
\end{corollary}

For example, \(z\mapsto z^k\) on \(S^1\) is a homotopy equivalence exactly when
\(k=\pm1\).

\section{Regular-value computation protocol}

For a smooth map \(f:M^n\to N^n\), a practical degree computation is:

\begin{enumerate}
\item choose a regular value \(y\in N\);
\item solve \(f(x)=y\);
\item determine whether each \(df_x\) preserves or reverses orientation;
\item sum the local signs.
\end{enumerate}

\begin{warning}[Unsigned counting is not degree]
\label{warn:viii03-unsigned}
The number \(|f^{-1}(y)|\) is generally not homotopy invariant.  Oppositely
oriented sheets can be created or cancelled in pairs.  Degree is the
\emph{signed} count.
\end{warning}

\section{Exercises}

\begin{exercise}\label{exr:viii03-01}
Define the local sign of a smooth map at a regular point.
\end{exercise}
\begin{hint}Inspect whether the derivative preserves orientation.\end{hint}
\begin{solution}
It is \(+1\) if \(df_x\) preserves orientation and \(-1\) if it reverses orientation.
\end{solution}

\begin{exercise}\label{exr:viii03-02}
Write the regular-value formula for degree.
\end{exercise}
\begin{hint}Sum local signs over the inverse image.\end{hint}
\begin{solution}
\[
\deg(f)=\sum_{x\in f^{-1}(y)}\operatorname{sgn}_x(f).
\]
\end{solution}

\begin{exercise}\label{exr:viii03-03}
What is the degree of the identity map?
\end{exercise}
\begin{hint}A regular value has one preimage.\end{hint}
\begin{solution}
\(1\).
\end{solution}

\begin{exercise}\label{exr:viii03-04}
What is the degree of a constant map \(S^n\to S^n\) for \(n>0\)?
\end{exercise}
\begin{hint}Choose a point not equal to the constant value.\end{hint}
\begin{solution}
\(0\).
\end{solution}

\begin{exercise}\label{exr:viii03-05}
If \(f\simeq g\), how are their degrees related?
\end{exercise}
\begin{hint}Degree is a homotopy invariant.\end{hint}
\begin{solution}
\(\deg(f)=\deg(g)\).
\end{solution}

\begin{exercise}\label{exr:viii03-06}
Why can a nonzero-degree map not be null-homotopic?
\end{exercise}
\begin{hint}Compare with a constant map.\end{hint}
\begin{solution}
A null-homotopic map has the degree of a constant map, namely \(0\).
\end{solution}

\begin{exercise}\label{exr:viii03-07}
State the multiplicativity formula.
\end{exercise}
\begin{hint}Take the degree of a composite.\end{hint}
\begin{solution}
\(\deg(g\circ f)=\deg(g)\deg(f)\).
\end{solution}

\begin{exercise}\label{exr:viii03-08}
What degrees are possible for a homotopy equivalence of oriented closed manifolds?
\end{exercise}
\begin{hint}Its degree must be a unit in \(\mathbb Z\).\end{hint}
\begin{solution}
Only \(+1\) and \(-1\).
\end{solution}

\begin{exercise}\label{exr:viii03-09}
Compute \(\deg(z\mapsto z^7)\) on \(S^1\).
\end{exercise}
\begin{hint}Use the power-map theorem.\end{hint}
\begin{solution}
\(7\).
\end{solution}

\begin{exercise}\label{exr:viii03-10}
Compute \(\deg(z\mapsto z^{-4})\) on \(S^1\).
\end{exercise}
\begin{hint}Negative exponent reverses winding orientation.\end{hint}
\begin{solution}
\(-4\).
\end{solution}

\begin{exercise}\label{exr:viii03-11}
When are two maps \(S^1\to S^1\) homotopic?
\end{exercise}
\begin{hint}Degree is complete in this case.\end{hint}
\begin{solution}
Exactly when they have the same degree.
\end{solution}

\begin{exercise}\label{exr:viii03-12}
What is the degree of a reflection of \(S^n\)?
\end{exercise}
\begin{hint}A reflection reverses orientation.\end{hint}
\begin{solution}
\(-1\).
\end{solution}

\begin{exercise}\label{exr:viii03-13}
Compute the degree of the antipodal map on \(S^1\).
\end{exercise}
\begin{hint}Use \((-1)^{n+1}\) with \(n=1\).\end{hint}
\begin{solution}
\(+1\).
\end{solution}

\begin{exercise}\label{exr:viii03-14}
Compute the degree of the antipodal map on \(S^2\).
\end{exercise}
\begin{hint}Use \((-1)^{n+1}\) with \(n=2\).\end{hint}
\begin{solution}
\(-1\).
\end{solution}

\begin{exercise}\label{exr:viii03-15}
Compute the degree of the antipodal map on \(S^3\).
\end{exercise}
\begin{hint}Use parity.\end{hint}
\begin{solution}
\(+1\).
\end{solution}

\begin{exercise}\label{exr:viii03-16}
If \(\deg f=2\) and \(\deg g=-3\), compute \(\deg(g\circ f)\).
\end{exercise}
\begin{hint}Multiply.\end{hint}
\begin{solution}
\(-6\).
\end{solution}

\begin{exercise}\label{exr:viii03-17}
Can a degree-\(2\) self-map of \(S^n\) be a homotopy equivalence?
\end{exercise}
\begin{hint}A homotopy equivalence has degree \(\pm1\).\end{hint}
\begin{solution}
No.
\end{solution}

\begin{exercise}\label{exr:viii03-18}
Why does nonzero degree force surjectivity?
\end{exercise}
\begin{hint}A missed regular value has no preimages.\end{hint}
\begin{solution}
If the map missed a regular value, the signed inverse-image sum would be empty,
forcing degree \(0\).
\end{solution}

\begin{exercise}\label{exr:viii03-19}
Why is ordinary preimage cardinality not a homotopy invariant?
\end{exercise}
\begin{hint}Pairs with opposite local signs can appear or disappear.\end{hint}
\begin{solution}
Homotopies can create or cancel pairs of preimages whose local signs sum to
zero, changing cardinality while preserving degree.
\end{solution}

\begin{exercise}\label{exr:viii03-20}
What degree would a retraction argument force on
\(r\circ i:S^n\to S^n\)?
\end{exercise}
\begin{hint}A retraction satisfies \(ri=\operatorname{id}\).\end{hint}
\begin{solution}
It would have degree \(1\).
\end{solution}

\begin{exercise}\label{exr:viii03-21}
Why is \(r\circ i\) null-homotopic if \(i:S^n\hookrightarrow D^{n+1}\)?
\end{exercise}
\begin{hint}The disk is contractible.\end{hint}
\begin{solution}
The inclusion \(i\) is homotopic in the disk to a constant map; postcomposing
that homotopy by \(r\) makes \(ri\) null-homotopic.
\end{solution}

\begin{exercise}\label{exr:viii03-22}
What contradiction proves the no-retraction theorem?
\end{exercise}
\begin{hint}Compare degrees \(1\) and \(0\).\end{hint}
\begin{solution}
A retraction would make \(ri=\operatorname{id}_{S^n}\) have degree \(1\), while
the contraction through the disk would make \(ri\) null-homotopic and hence
degree \(0\).
\end{solution}

\begin{exercise}\label{exr:viii03-23}
If \(\deg f=1\), must \(f\) be a homeomorphism?
\end{exercise}
\begin{hint}Degree measures global signed covering, not injectivity.\end{hint}
\begin{solution}
No.  Degree \(1\) does not imply injectivity or local homeomorphism behavior.
\end{solution}

\begin{exercise}\label{exr:viii03-24}
What topic does the antipodal degree formula prepare for next?
\end{exercise}
\begin{hint}Look at VIII/04.\end{hint}
\begin{solution}
It prepares the study of spheres and antipodal maps, including parity-sensitive
obstructions and Borsuk--Ulam-type phenomena.
\end{solution}

\section{Solved problem dossiers}

\begin{problem}[Power map]\label{prob:viii03-01}
Compute the degree of \(p_6(z)=z^6\).
\end{problem}
\begin{solution}
A regular value has six preimages, each orientation preserving, so
\(\deg(p_6)=6\).
\end{solution}

\begin{problem}[Negative power]\label{prob:viii03-02}
Compute the degree of \(q(z)=z^{-5}\).
\end{problem}
\begin{solution}
The map winds five times with reversed orientation, so \(\deg(q)=-5\).
\end{solution}

\begin{problem}[Composition]\label{prob:viii03-03}
Let \(f(z)=z^3\) and \(g(z)=z^{-2}\).  Compute the degree of \(g\circ f\).
\end{problem}
\begin{solution}
\[
\deg(g\circ f)=(-2)(3)=-6.
\]
Indeed \(g(f(z))=z^{-6}\).
\end{solution}

\begin{problem}[Homotopy obstruction]\label{prob:viii03-04}
Can \(z\mapsto z^2\) be homotopic to \(z\mapsto z^3\)?
\end{problem}
\begin{solution}
No.  Their degrees are \(2\) and \(3\), and homotopic maps have equal degree.
\end{solution}

\begin{problem}[Null-homotopy test]\label{prob:viii03-05}
Can \(z\mapsto z^{-1}\) be null-homotopic as a map \(S^1\to S^1\)?
\end{problem}
\begin{solution}
No.  Its degree is \(-1\), whereas a null-homotopic circle map has degree \(0\).
\end{solution}

\begin{problem}[Antipodal parity]\label{prob:viii03-06}
For which \(n\) does the antipodal map \(S^n\to S^n\) have degree \(+1\)?
\end{problem}
\begin{solution}
When \(n\) is odd, because \((-1)^{n+1}=+1\).
\end{solution}

\begin{problem}[Antipodal parity II]\label{prob:viii03-07}
For which \(n\) does the antipodal map have degree \(-1\)?
\end{problem}
\begin{solution}
When \(n\) is even.
\end{solution}

\begin{problem}[No retraction]\label{prob:viii03-08}
Give the degree contradiction for a hypothetical retraction
\(D^2\to S^1\).
\end{problem}
\begin{solution}
For inclusion \(i:S^1\hookrightarrow D^2\), a retraction \(r\) gives
\(ri=\operatorname{id}_{S^1}\), degree \(1\).  But \(i\) contracts through
\(D^2\), so \(ri\) is null-homotopic, degree \(0\).  Contradiction.
\end{solution}

\begin{problem}[Surjectivity]\label{prob:viii03-09}
Suppose \(f:S^2\to S^2\) has degree \(4\).  What can be said about its image?
\end{problem}
\begin{solution}
It must be all of \(S^2\), since nonzero degree implies surjectivity.
\end{solution}

\begin{problem}[Homotopy equivalence test]\label{prob:viii03-10}
Suppose \(f:S^3\to S^3\) has degree \(-1\).  Does degree alone prove that \(f\)
is a homotopy equivalence?
\end{problem}
\begin{solution}
Degree \(-1\) is necessary and, for sphere self-maps, it is sufficient because
sphere maps of the same degree are homotopic and a degree \(-1\) reflection is
a homeomorphism.  In general manifolds, degree \(\pm1\) alone need not imply
homotopy equivalence.
\end{solution}

\begin{problem}[Classification on the circle]\label{prob:viii03-11}
Classify the homotopy class of \(f(z)=\overline z^{\,4}\).
\end{problem}
\begin{solution}
Since \(\overline z=z^{-1}\) on \(S^1\),
\(f(z)=z^{-4}\), so \(\deg f=-4\).  Its homotopy class is the unique circle-map
class of degree \(-4\).
\end{solution}

\begin{problem}[Bridge to VIII/04]\label{prob:viii03-12}
Explain what the formula \(\deg(a)=(-1)^{n+1}\) says qualitatively about
antipodal maps.
\end{problem}
\begin{solution}
Their homotopy behavior depends on the parity of the sphere dimension:
odd-dimensional spheres give degree \(+1\), even-dimensional spheres give
degree \(-1\).  VIII/04 turns this parity into concrete sphere and antipodal-map
phenomena.
\end{solution}
